% TODOS: Citations (ctrl+f TODO) and Bib file
% PaddleOCR-V-1.6 links: https://huggingface.co/PaddlePaddle/PaddleOCR-VL-1.6, technical report https://arxiv.org/pdf/2606.03264, bibtex
% OpenKB: https://github.com/VectifyAI/OpenKB
% PageIndex: https://github.com/VectifyAI/PageIndex, bibtex

\documentclass[11pt, a4paper]{article}
\usepackage[margin=2.3cm]{geometry}
\usepackage{parskip}
\usepackage{enumitem}
\usepackage{hyperref}
% biblatex with the bibtex backend (no biber dependency, so it compiles under
% tectonic).  Switch to backend=biber on a full TeX Live install if preferred.
\usepackage[backend=bibtex,style=ieee]{biblatex}
\addbibresource{references.bib}


\title{\textbf{Structure-Faithful PDF Ingestion for Tree-Based RAG}}
\author{Federico Mizzaro, Omar Chafik, Abdul Rahim}

\begin{document}
\maketitle

% ==============================================================
\section*{Motivation}
% ==============================================================

Hospital infection-prevention and control (IPC) teams rely on
multi-level clinical guidelines to make (time-sensitive) decisions.
Recent vectorless RAG frameworks such as PageIndex\cite{zhang} would be
able to make use of such document hierarchy for reasoning-based,
top-down retrieval boasting extreme accuracy\cite{benchmark}.  However,
these tools assume clean, correctly-structured input of the document
set: the heading hierarchy and leaf content must survive ingestion
intact. The PageIndex documentation itself warns that for markdown
converted from PDF, the de facto format for published and archived
documents, ``most existing conversion tools cannot preserve the original
hierarchy'' \cite{pageindexrepo}.

The PageIndex ecosystem's own answer to this problem, used by the OpenKB
knowledge base\cite{openkb}, converts documents with Microsoft's
\texttt{markitdown} and extracts images with \texttt{pymupdf}. For
``large or complex PDFs,'' both projects direct users to the
proprietary, hosted \emph{PageIndex Cloud} OCR service. No open-source,
locally-deployable OCR ingestion path exists yet in this ecosystem.

For clinical deployment, a cloud OCR dependency is questionable at best,
and the quality gap left by text-layer extraction alone is a
patient-safety risk. Our target deployment is therefore deliberately
constrained: \textbf{commodity, CPU-only, fully local} --- no GPU, no
hosted service. This project addresses the gap directly: we build an
open ingestion path on top of the PaddleOCR-VL-1.6 \cite{paddle} model
family, and we supply the reproducible fidelity benchmark that the
ecosystem currently lacks to measure the quality of such a path.

% ==============================================================
\section*{Analysis}
% ==============================================================

The project has two separable contributions.

\textbf{Contribution A: a benchmark harness.} We construct a
ground-truth corpus of structured PDFs by compiling LaTeX documents to
PDF while retaining the source. Ground-truth heading trees are extracted
directly from \texttt{\textbackslash section}, \texttt{\textbackslash
subsection} (and so on) commands, and ground-truth asset inventories
from (as examples) \texttt{\textbackslash includegraphics} and
\texttt{\textbackslash begin\{table\}} declarations. This gives us
unambiguous, machine-readable ground truth without human annotation. The
harness then runs competing ingestion paths --- the OpenKB-style
\texttt{markitdown}+LLM reconstruction (the current default) and our OCR
pipeline --- and compares the resulting trees against ground truth. A
central methodological commitment is that we score \emph{detection} and
\emph{nesting} \textbf{separately} (see Implementation), so a win in
finding the right headings is never confused with a win in nesting them
correctly.

\textbf{Contribution B: an open, CPU-tractable OCR ingestion pipeline.}
The naive use of PaddleOCR-VL --- running the 0.9B vision-language model
over every block --- is far too slow to be local and practical (a single
14-page paper took $\sim$48\,min on our CPU-only target). Our pipeline
instead recovers the heading \emph{tree} with a fast detection model and
a lightweight recogniser, reserving the heavy VLM for optional content
recognition. On top of detection we add a \emph{leveling} layer that
turns flat heading blocks into a relative hierarchy, driven primarily by
the headings' own rendered numbering and falling back to a font-size
signal when numbering is absent. A planned third stage (Contribution
B$'$, below) introduces an \emph{anchored-asset leaf} scheme that
attaches figures and tables as dedicated, addressable leaf nodes in the
PageIndex tree rather than inline blobs the LLM happens to read.

% ==============================================================
\section*{Implementation}
% ==============================================================

Main libraries: \texttt{paddlex}/\texttt{paddleocr} (PP-DocLayoutV3 +
mobile PP-OCR), \texttt{pypdfium2}, \texttt{markitdown}, and the
open-source PageIndex codebase.

\textbf{Document layout analysis (the speed decision).} We found
empirically that the VLM's cost is dominated by \emph{autoregressive text
generation} over every block, not vision encoding --- shrinking the input
image made it \emph{slower}, and no single knob (\texttt{max\_pixels},
\texttt{max\_new\_tokens}, disabling sub-pipelines) made it tractable.
The win is architectural: for the heading tree, only the short
\emph{title} blocks need to be read. We therefore use
\texttt{PP-DocLayoutV3}, a detection model with \emph{no} generation, to
label every block (\texttt{doc\_title}, \texttt{paragraph\_title},
\texttt{text}, \texttt{figure\_title}, \texttt{table}, \ldots) with a
bounding box at $\sim$1--4\,s/page, and a mobile \texttt{PaddleOCR}
(classic lite det+rec) to recognise only the title crops --- at
\emph{identical} title quality to the VLM. This is one to two orders of
magnitude faster than the full VLM, which remains available for the
separate, slower content-fidelity pass.

\textbf{Heading detection, not heading levels.} The layout model gives a
reliable \emph{label} but not a reliable level. We isolate real section
headings by label alone (\texttt{paragraph\_title}), with \emph{no} text
heuristics: captions carry a distinct \texttt{figure\_title} label, so
filtering on the label --- rather than a rule like ``drop titles starting
with \emph{Table}'' --- avoids overfitting the evaluator to any one
document. (A rare caption mislabelled \texttt{paragraph\_title} is then
an honest false positive, not a hidden patch.)

\textbf{Leveling.} Levels are reconstructed from the signal already
present in the recognised title text: its \textbf{rendered numbering}.
Dotted arabic (\texttt{1}, \texttt{1.1}, \texttt{1.2.3}) gives level by
component count; IEEE-style \texttt{I.}/\texttt{A.}/\texttt{1)} maps to
section/subsection/subsubsection. The hard case --- single tokens like
\texttt{I}, \texttt{V}, \texttt{C}, \texttt{D} that are valid roman
numerals \emph{and} subsection letters, made worse because IEEEtran
renders sections in Title Case (so capitalisation is uninformative) ---
is resolved by \textbf{sequence}: a token is a section if it continues
the roman run, a subsection if it continues the letter run. For
\emph{unnumbered} documents we fall back to an independent
\textbf{text-layer font-size} signal read via \texttt{pypdfium2} (largest
size = level~1); font size is unreliable in general (it inverts on IEEE
small-caps) but is safe here precisely because numbered docs take the
numbering path, so the font fallback only fires where size decreases
cleanly with depth. Throughout, we interpret the heading sequence as a
\emph{relative} hierarchy: a level gap (a heading nested more than one
step below its parent, possibly an authoring artefact) is absorbed as the
next nesting step, not treated as an error to repair. A column-aware
\textbf{reading-order} sort is applied first, since two-column layouts
otherwise return blocks out of order (e.g.\ placing \texttt{2.1} before
section~\texttt{2}).

\textbf{Baseline.} The comparison point is the OpenKB/PageIndex path:
\texttt{markitdown} dumps the PDF text layer to flat markdown (no
headings) and an LLM reconstructs the heading list. We make this a
\emph{strong} and \emph{deterministic} baseline --- the reconstruction is
produced by a capable LLM \textbf{blind to the LaTeX ground truth},
cached as one frozen artefact per document, so runs are reproducible and
we are beating the conservative target rather than a weak local model.

\textbf{Benchmark computation.} LaTeX source is parsed to ground-truth
trees and asset inventories. We replace a single tree-edit-distance score
with two decomposable F1 metrics over normalised titles: \texttt{title\_f1},
the level-free bag overlap of section titles (pure \emph{detection}), and
\texttt{edge\_f1}, the overlap of parent$\to$child pairs (\emph{nesting}
and identity together). Both use a \texttt{normalise\_title} step that
strips rendered numbering and markdown emphasis, and both are computed
identically for every condition. Content fidelity (recognising body prose
under each heading) is a separate, slower concern handled by the full VLM
pass and scored by token-level edit distance over aligned leaves; the
fast structure path deliberately does not attempt it. OCR results are
cached by \texttt{(pdf, config)}; we verified the OCR output is
bit-stable across repeated runs, so the benchmark is fully reproducible.
The ingest module slots into the existing PageIndex
\texttt{modules/ingest/} registry, altering only how the tree is built
and leaving retrieval untouched: the goal is not to improve retrieval but
to make the OCR path open, local, and measurable.

\textbf{Contribution B$'$ (planned): anchored assets.} Figures and tables
will be attached as dedicated leaf nodes anchored to the prose leaf under
which they appear (by spatial association from the layout bounding
boxes), making non-text content a first-class, addressable part of the
retrieval tree. This stage is designed but not yet implemented; the
evaluation below treats it as the final ablation step.

% ==============================================================
\section*{Evaluation}
% ==============================================================

We evaluate ingestion fidelity comparatively, against the OpenKB
\texttt{markitdown}+LLM reconstruction baseline, on a LaTeX ground-truth
corpus spanning the layout and numbering axes: synthetic documents in
distinct templates (single- and two-column \texttt{article} with arabic
numbering, \texttt{IEEEtran} with roman/letter/paren numbering, and
unnumbered \texttt{article} both flat and three-level nested), plus real
two-column IEEE papers --- seven documents in total, each compiled from
source with \texttt{tectonic}. The design is a four-condition ablation
that isolates each component: (A) \texttt{markitdown}+LLM (OpenKB
default); (B) OCR title detection with \emph{no} leveling (flat tree);
(C) OCR with numbering/font leveling; (D) the full pipeline including
anchored-asset leaves. Detection is scored by \texttt{title\_f1}
(A vs.\ B), nesting by \texttt{edge\_f1} (C), and the asset contribution
by the (C)$\to$(D) contrast, measured as asset recall: the fraction of
\texttt{\textbackslash includegraphics} / \texttt{\textbackslash
begin\{table\}} declarations recovered as anchored leaves.

This design is suitable because it is fully automatable, needs no
subjective human annotation, compares \emph{text against text} at every
stage (avoiding the ambiguity of PDF-versus-text comparison), and
directly measures the gap between the current open-source default and our
pipeline --- attributing each gain to a specific component rather than to
the pipeline as a whole.

Early results on the assembled corpus already support the core thesis.
On a hard real two-column IEEE paper, OCR recovers \emph{every} heading
(\texttt{title\_f1} $\approx 0.97$, recall $1.0$) where the
\texttt{markitdown} baseline reaches only $\approx 0.64$ (recall $0.53$):
text-layer extraction interleaves the two columns and even drops
inter-word spaces, scattering the \texttt{1)}/\texttt{2)} subsubsections
beyond recovery, whereas reading the rendered image per block isolates
each heading regardless of column flow. On a clean single-column document
the two tie at $1.0$ --- OCR earns its cost on hard layouts, not on easy
ones, which we state plainly. With leveling enabled, five of the seven
documents reconstruct the tree \emph{exactly} (\texttt{edge\_f1} $= 1.0$).
The two residual cases are \textbf{characterised, not patched}: IEEEtran
renders \texttt{\textbackslash subsubsection} run-in (inline, so those
headings do not exist as blocks in the rendered PDF --- a rendering
boundary OCR cannot cross), and one table caption mislabelled by the
layout model cascades into its neighbours' nesting. We deliberately do
not add document-specific rules to erase these, so the reported numbers
reflect the pipeline's true behaviour rather than a tuned evaluator.

% ==============================================================


\printbibliography


\end{document}
